% ============================================================
%  Section 4 — Experimental Setup
%  Dataset-specific choices and the evaluation protocol. The
%  framework itself is defined in Section 3; measured outcomes are
%  in Section 5.
% ============================================================

\section{Experimental Setup}
\label{sec:experimental_setup}

% ------------------------------------------------------------
\subsection{Dataset Description}
\label{sec:dataset}

The framework is evaluated on the XJTU-SY rolling element bearing
dataset~\cite{8576668, lei2019xjtusy}, a widely used benchmark that
provides complete run-to-failure vibration data from accelerated life
tests. The dataset comprises 15~bearings in three groups of five, each
tested under a distinct operating condition
(Table~\ref{tab:xjtu_conditions}). Two accelerometers record the
horizontal and vertical vibration simultaneously. Each recording lasts
1.28~s (32,768~samples at 25.6~kHz) and is taken once per minute until
failure.

\begin{table}[htbp]
\centering
\caption{XJTU-SY operating conditions and bearing characteristics.}
\label{tab:xjtu_conditions}
\footnotesize
\renewcommand{\arraystretch}{1.2}
\begin{tabular}{cccccc}
\toprule
\textbf{Cond.} & \textbf{Speed} & \textbf{Load} &
\textbf{Bearings} & \textbf{Lifetime range} &
\textbf{Failure modes} \\
& (Hz) & (kN) & & (min) & \\
\midrule
1 & 35.0  & 12.0 & 1-1 to 1-5 & 52--161 &
  OR (3), Cage (1), OR+Ball (1) \\
2 & 37.5  & 11.0 & 2-1 to 2-5 & 42--533 &
  OR (3), IR (1), Cage (1) \\
3 & 40.0  & 10.0 & 3-1 to 3-5 & 114--2538 &
  OR (2), IR (2), IR+Ball (1) \\
\bottomrule
\end{tabular}

\vspace{2pt}
\raggedright\scriptsize OR: outer race. IR: inner race.
\end{table}

Bearing lifetimes vary widely across and within operating conditions
(Fig.~\ref{fig:lifetime_distribution}). Condition~1, with the heaviest
radial load, produces the shortest lifetimes (52--161~min), while
Condition~3 ranges from 114 to over 2,500~min. Within Condition~2,
Bearing~2-4 failed after 42~min and Bearing~2-3 after 533~min under
identical settings. Failure modes also differ between bearings of the
same condition, so the operating parameters influence the failure mode
without determining it.\\

\begin{figure}[htbp]
  \centering
  \includegraphics[width=\columnwidth]{figures/01_lifetime_distribution.png}
  \caption{Bearing lifetimes across the three operating conditions
  of the XJTU-SY dataset.}
  \label{fig:lifetime_distribution}
\end{figure}

Fig.~\ref{fig:all_bearings_rms} shows the RMS of the horizontal
channel over the life of all 15~bearings. The RMS stays near a low
baseline during healthy operation, rises as degradation starts, and
climbs sharply before failure. The dynamics nonetheless differ between
bearings: smooth monotone growth (Bearing~3-1), abrupt failure with
little warning (Bearing~1-4, cage fault), oscillation with partial
recoveries (Bearing~2-5), or several distinct rates of change
(Bearing~1-2). A predictor must learn from all of these patterns, and
an interval must remain valid across them.

\begin{figure}[htbp]
  \centering
  \includegraphics[width=\columnwidth]{figures/01_all_15_bearings_rms_overview.png}
  \caption{RMS vibration amplitude (horizontal channel) over the
  operational life of all 15~XJTU-SY bearings, organised by operating
  condition, with bearing ID, failure mode and lifetime.}
  \label{fig:all_bearings_rms}
\end{figure}

% ------------------------------------------------------------
\subsection{Health Indicator, Stage Labels and Features}
\label{sec:hi_staging}

The health indicator of Section~\ref{sec:method_staging} is the RMS of
the horizontal channel, which responds most strongly to the radial
load applied in these tests~\cite{lu_joint_2022}, smoothed with a
short moving average. The healthy baseline is estimated over the first
$\eta = 20\%$ of each bearing's life, the relative floor is
$\rho = 0.20$, and the persistence length is $m = 5$ recordings.\\

The relative floor is necessary on this dataset. Long-lived bearings
have an exceptionally stable baseline: for Bearing~3-1 the healthy RMS
varies by only a few percent, so $\mu + 3\sigma$ lies barely above the
healthy level and is crossed by ordinary fluctuations. Without the
floor, the FPT of Bearing~3-1 is detected at minute~23 of 2,537,
labelling more than 99\% of its life as degrading, and that of
Bearing~3-2 at minute~169 of 2,495. These two bearings hold over half
of all recordings, so the error would propagate into every partition.
Across the 15~bearings, healthy fluctuations remain within about 10\%
of the baseline, whereas degradation raises the RMS by close to an
order of magnitude, so a 20\% floor separates the two regimes with a
wide margin.\\

Fig.~\ref{fig:hi_all_bearings} overlays the resulting stages on the
health indicator. In every bearing the onset of Stage~2 coincides with
the visible departure from the baseline and Stage~3 with the terminal
rise. Two groups of bearings appear. Some degrade abruptly, with a flat
indicator for most of their life and a near-vertical rise that
compresses both degradation stages into a short final window
(Bearings~2-1, 3-1, 3-3 and~3-4). Others degrade progressively, often
with oscillations (Bearings~1-2, 2-3, 2-5 and~3-2). Abrupt failures
occur in short- and long-lived bearings alike, which is why the
threshold is expressed in relative terms.

\begin{figure}[htbp]
  \centering
  \includegraphics[width=\textwidth]{figures/02_all_bearings_stages.png}
  \caption{Health indicator (horizontal RMS) of all 15~bearings with
  the automatically detected stages. Green: Stage~1 (healthy). Orange:
  Stage~2 (early degradation). Red: Stage~3 (near-failure). Dashed
  lines: FPT and acceleration point.}
  \label{fig:hi_all_bearings}
\end{figure}

Over all 9,216~recordings, Stage~1 accounts for 70.7\%, Stage~2 for
15.9\% and Stage~3 for only 13.4\% (Fig.~\ref{fig:stage_distribution}).
The imbalance is stronger still within bearings: several bearings
contribute only a handful of near-failure recordings, and Bearing~3-2
alone holds about half of all near-failure windows. This is the
scarcity that C1 targets, and the concentration that makes calibration
on few bearings fragile.

\begin{figure}[htbp]
  \centering
  \includegraphics[width=\columnwidth]{figures/02_stage_distribution_per_bearing_pct.png}
  \caption{Degradation stage distribution. Left: composition of each
  bearing's life. Right: global composition over all 9,216~recordings.}
  \label{fig:stage_distribution}
\end{figure}

The predictors use a multivariate representation. Ten time-domain
features are extracted per channel and recording: amplitude
descriptors (RMS, peak, peak-to-peak, standard deviation),
distribution-shape descriptors (kurtosis, skewness) and dimensionless
impulse descriptors (crest, shape, impulse and clearance factors).
The two channels give $d = 20$ features per recording.
Fig.~\ref{fig:feature_matrix} shows that these features are stable
during the healthy phase and change sharply from the FPT onward, with
amplitude and distribution-shape descriptors moving in opposite
directions. The feature families therefore carry complementary
information, which motivates generating and predicting from the full
multivariate window rather than a scalar indicator.

\begin{figure}[htbp]
  \centering
  \includegraphics[width=\columnwidth]{figures/02_feature_matrix_heatmap.png}
  \caption{Temporal evolution of the 20 time-domain features (z-scored
  per feature) for Bearing~2-1. Dashed lines: FPT and acceleration
  point.}
  \label{fig:feature_matrix}
\end{figure}

% ------------------------------------------------------------
\subsection{Windowing, Normalisation and Partitioning}
\label{sec:windowing}

Windows of $w = 32$ recordings are taken with a stride of one, giving
inputs of dimension $32 \times 20$. The unit stride is kept despite the
overlap between adjacent windows because the shortest bearings would
otherwise contribute almost nothing: Bearing~2-4 yields only
11~windows even at unit stride.\\

The amplitude, impulse and kurtosis features are strictly positive and
grow by about a factor of twenty near failure. They are log-transformed
before standardisation, which would otherwise be dominated by the
abundant healthy windows and place near-failure values tens of
standard deviations from the mean. Standardisation statistics are
estimated on the training windows of each model only, and the
standardised values are clipped to $\pm 5$ standard deviations to bound
the low-variance shape features.\\

\textit{Outer folds.} Five folds are formed following the protocol of
Lu et al.~\cite{lu_joint_2022}: fold~$k$ tests on Bearings~1-$k$, 2-$k$
and~3-$k$, so that every fold spans the three operating conditions and
every bearing is tested exactly once.\\

\textit{Inner groups for cross-conformal calibration.} The 12 non-test
bearings of each fold contain four bearings per operating condition.
They are divided into $K = 4$ groups of three, group~$j$ holding the
$j$-th remaining bearing of each condition, so that every inner model
is trained on all three conditions, as the test set is. Each inner
model therefore trains on 9~bearings and scores its 3~held-out
bearings. Of the 9, one is used for early stopping: the smallest
bearing that yields at least 300~windows, so that the validation error
is stable enough to detect a genuine plateau, while as few windows as
possible are withheld from training. The remaining 8~bearings form the
training set. Over the 12 non-test bearings, the four inner models
provide out-of-sample residuals for every window.\\

Because three bearings account for most of the recordings, the folds
remain unbalanced after partitioning: test sets range from 412 to
3,059~windows. Per-fold metrics are therefore estimated with unequal
precision, which is why results are also examined per bearing
(Section~\ref{sec:res_per_bearing}).

% ------------------------------------------------------------
\subsection{Implementation Details}
\label{sec:implementation}

Table~\ref{tab:hyperparams} lists the hyperparameters. The augmentation
ratio $r = 0.75$ and the stage-loss weight $\lambda = 1.0$ were
selected on preliminary runs over $r \in \{0.25, 0.5, 0.75, 1.0\}$ and
$\lambda \in \{0.3, 0.5, 1.0\}$. A generator is trained for each of the
20 inner models (5 folds $\times$ 4 groups); it does not depend on the
seed. In the single-task configurations, the stage used by C3 comes
from a stage network identical in architecture and training to the
multi-task predictor without augmentation.

\begin{table}[htbp]
\centering
\caption{Hyperparameters of the framework components.}
\label{tab:hyperparams}
\footnotesize
\renewcommand{\arraystretch}{1.15}
\begin{tabular}{lll}
\toprule
\textbf{Component} & \textbf{Setting} & \textbf{Value} \\
\midrule
\multirow{8}{*}{C1 (generator, critic)}
 & Noise dimension & 64 \\
 & Conv.\ channels / kernel & 128 / 5 \\
 & BiLSTM units per direction & 128 \\
 & Optimiser, learning rate & Adam, $10^{-4}$ \\
 & Adam $(\beta_1, \beta_2)$ & $(0.5, 0.9)$ \\
 & Batch size / epochs & 64 / 300 \\
 & Critic steps $n_{\mathrm{c}}$ & 3 \\
 & Penalty weight $\lambda_{\mathrm{gp}}$ & 10 \\
\midrule
\multirow{8}{*}{C2 (predictor)}
 & BiLSTM layers / units & 2 / 48 \\
 & Head hidden units & 48 \\
 & Dropout & 0.4 \\
 & Optimiser, learning rate & Adam, $5 \times 10^{-4}$ \\
 & Weight decay & $5 \times 10^{-4}$ \\
 & Batch / max.\ epochs / patience & 64 / 100 / 25 \\
 & Gradient clipping (norm) & 5.0 \\
 & Stage-loss weight $\lambda$ & 1.0 \\
\midrule
Augmentation & Ratio $r$ & 0.75 \\
\multirow{2}{*}{C3 (conformal)} & Inner groups $K$ & 4 \\
 & Miscoverage $\alpha$ & 0.10\textsuperscript{a} \\
\bottomrule
\end{tabular}

\vspace{2pt}
\raggedright\scriptsize \textsuperscript{a}Headline level; 0.05, 0.15
and 0.20 are also evaluated.
\end{table}

All models were implemented in PyTorch and trained on a single NVIDIA
T4 GPU. Recurrent training on the GPU is not bit-wise reproducible, and
the RUL error of a given configuration varies by about
$\pm 0.015$--$0.02$ from run to run. Every configuration is therefore
trained under five seeds; with 5~folds and 4~inner models, each
configuration comprises 100 trained predictors. Results are the mean
and standard deviation over seeds of the fold-averaged metrics, and a
difference smaller than the seed-to-seed variability is not interpreted
as an effect.

% ------------------------------------------------------------
\subsection{Evaluation Metrics}
\label{sec:metrics}

Point accuracy is the root mean square error on the normalised RUL
scale, computed per test bearing and then averaged, so that the
long-lived bearings, which supply most windows, do not dominate:
\begin{equation}
  \mathrm{RMSE} = \frac{1}{|\mathcal{B}_{\mathrm{te}}|}
  \sum_{b \in \mathcal{B}_{\mathrm{te}}}
  \sqrt{ \frac{1}{|\mathcal{W}_b|} \sum_{t \in \mathcal{W}_b}
  \left( \hat{y}^{b}_{t} - y^{b}_{t} \right)^2 } ,
  \label{eq:rmse}
\end{equation}
where $\mathcal{W}_b$ is the set of evaluated windows of bearing $b$.
Interval quality is measured by the prediction interval coverage
probability and the mean prediction interval width,
\begin{equation}
  \mathrm{PICP} = \frac{1}{n} \sum_{i=1}^{n}
  \mathbb{1}\left\{ y_i \in \hat{C}(\mathbf{X}_i) \right\},
  \qquad
  \mathrm{MPIW} = \frac{1}{n} \sum_{i=1}^{n}
  \left| \hat{C}(\mathbf{X}_i) \right| ,
  \label{eq:picp_mpiw}
\end{equation}
computed over the test windows of each fold, with intervals clipped to
$[0,1]$. PICP should lie at or slightly above $1-\alpha$, and a smaller
MPIW is better only when PICP is maintained, so the two are always
reported together. All metrics are reported for the healthy, early and
near-failure stages, for the post-FPT windows and for all windows.
Stage subsets are always defined by the true stage, so that the
configurations are compared on the same windows.\\

Three versions of C3 are compared, all built on the same CV+
residuals: a single quantile for all windows; quantiles per
\emph{predicted} stage, which is the deployable framework; and
quantiles per \emph{true} stage, which uses information available only
after failure and serves as a reference for what perfect stage
recognition would achieve. The stage prediction itself is assessed by
its accuracy and by the recall and precision of the near-failure
stage.
